\documentclass[11pt]{article}

\usepackage[margin=1in]{geometry}
\usepackage{amsmath, amssymb}
\usepackage{tabularx}
\usepackage{xcolor}
\usepackage[colorlinks=true, linkcolor=blue!50!black, citecolor=blue!50!black, urlcolor=blue!50!black]{hyperref}

\newcommand{\F}{\mathbb{F}}
\newcommand{\K}{K}            % F_{2^64}: addresses, counters, committed lanes
\newcommand{\E}{E}            % F_{2^192} = K[y]/(y^3+y+1): protocol/extension values
\newcommand{\pkd}{\mathrm{pkd}}
\newcommand{\eq}{\operatorname{eq}}
\newcommand{\cube}[1]{\{0,1\}^{#1}}

% VM macros
\newcommand{\gen}{\mathsf{g}}
\newcommand{\loc}[1]{\mem[\fp\cdot #1]}
\newcommand{\pc}{\mathbf{pc}}
\newcommand{\fp}{\mathbf{fp}}
\newcommand{\mem}{\mathbf{m}}
\newcommand{\addr}{\mathit{addr}}
\newcommand{\val}{\mathit{val}}
\newcommand{\cnt}{\mathit{count}}
\newcommand{\push}{\textsc{push}}
\newcommand{\pull}{\textsc{pull}}
\newcommand{\tup}[1]{\langle #1 \rangle}
\newcommand{\dsep}[1]{\mathsf{#1}}
\newcommand{\opc}[1]{\mathsf{op}_{\mathsf{#1}}}
\newcommand{\srcsel}[1]{\mathsf{src}(#1)}

\title{\bf leanVM-b}
\author{}
\date{}

\begin{document}
\maketitle


\tableofcontents
\newpage

\section{VM specification}\label{sec:vm}

\paragraph{Field.}
\[
  \K=\F_{2^{64}}=\F_2[x]/(x^{64}+x^4+x^3+x+1),
  \qquad
  \E=\K[y]/(y^3+y+1),
  \qquad |\E|=2^{192}.
\]
A \emph{memory word}, immediate, hash lane, address, program counter $\pc$, frame pointer $\fp$, and read counter is one element of $\K$. Base addition is bitwise \textsc{xor} and base multiplication is in $\K$. An extension value $v=v_0+v_1y+v_2y^2\in\E$ is stored as three consecutive $\K$ words; explicit extension instructions reassemble those coordinates inside their constraints. Protocol randomness is sampled from $\E$ (\S\ref{sec:prelim}), while every physical committed column remains $\K$-valued.

\paragraph{Addresses as powers of a generator.} Fix a generator $\gen$ of the multiplicative group $\K^{\times}$, of order $2^{64}-1$. Memory addresses and bytecode addresses (the program counter) are powers of $\gen$: the $i$-th address is the field element $\gen^{\,i}$. Its successor is one field multiplication: $\gen^{\,i+1}\;=\;\gen\cdot\gen^{\,i}$.

\paragraph{Memory.} Memory is read-only (write-once): an array of $2^{h}$ cells, each holding one $64$-bit $\K$ word, addressed by the $2^{h}$ powers $\gen^{0},\gen^{1},\dots,\gen^{2^{h}-1}$. The size exponent $h$ is fixed per execution, with $16\le h\le 32$. Each cell is set once; every access thereafter reads it. Thus an extension value beginning at address $a$ occupies $\mem[a],\mem[\gen a],\mem[\gen^2a]$.

\paragraph{Program.} The program is a fixed, public sequence of instructions, addressed like memory by powers of $\gen$: the program counter $\gen^{\,i}$ selects the $i$-th instruction. An instruction is an opcode together with its operands, described below.

\paragraph{Registers.} The machine state is two registers, the program counter $\pc$ and the frame pointer $\fp$, each a power of $\gen$. Every instruction except \texttt{JUMP} advances the state by $\pc\gets\gen\cdot\pc$ (fetching the next instruction) and leaves $\fp$ unchanged; \texttt{JUMP} may set both.

\paragraph{Operands and addressing.} An operand is either an \emph{immediate}, a literal field element used as given, or an \emph{fp-relative reference}: an offset $j$ into the current frame, encoded as the $\gen$-power $o=\gen^{\,j}$. With the frame based at $\fp=\gen^{\mathrm{base}}$, the reference names the cell
\[
  \mem[\fp\cdot o],\qquad \fp\cdot o\;=\;\gen^{\mathrm{base}}\cdot\gen^{\,j}\;=\;\gen^{\,\mathrm{base}+j}.
\]
Since addresses are exponents of $\gen$, the frame is the contiguous block $\gen^{\mathrm{base}},\gen^{\mathrm{base}+1},\dots$, and adding the offset $j$ to the base exponent is, in the field, a single multiplication $\fp\cdot o$.

\paragraph{Instruction set.} The machine implements nine instruction-table families; two have bytecode-bound narrow/scalar modes used by the 64-bit representation. Below, $o,o_A,\dots$ are fp-relative references and $k$ a base-field immediate.
\begin{center}
\begin{tabularx}{\textwidth}{@{}llX@{}}
\hline
Instruction & Operands & Semantics\\
\hline
\texttt{XOR}           & $[o_A,o_B,o_C]$ & $\loc{o_C}=\loc{o_A}+\loc{o_B}$ in $\K$\\
\texttt{MUL}           & $[o_A,o_B,o_C]$ & $\loc{o_C}=\loc{o_A}\cdot\loc{o_B}$ in $\K$\\
\texttt{XOR\_192}      & $[o_A,o_B,o_C]$ & three-word extension addition\\
\texttt{MUL\_192}      & $[o_A,o_B,o_C;m]$ & extension multiplication; $m=1$ embeds scalar $A$\\
\texttt{SET\_CONSTANT} & $[o,k]$         & $\loc{o}=k$ \quad($k\in\K$)\\
\texttt{DEREF}         & $[\alpha,\beta,\gamma;\,s]$ & $\mem[\loc{\alpha}\cdot\beta]=\srcsel{s}$\\
\texttt{DEREF\_WIDE}   & $[\alpha,\beta,\gamma;w]$ & equality of $2+w$ consecutive words\\
\texttt{JUMP}          & $[o_c,o_d,o_f]$ & conditional jump (below)\\
\texttt{BLAKE3}        & $[o_{A,0},o_{A,1},o_{B,0},o_{B,1},o_{CV},o_C;m_0,m_1]$ & $c=\mathrm{BLAKE3}(a,b)$ (below)\\
\hline
\end{tabularx}
\end{center}

For \texttt{XOR\_192} and \texttt{MUL\_192}, each operand reference names its first coordinate. If $a_A=\fp o_A$, the other two addresses are the virtual columns $\gen a_A$ and $\gen^2a_A$; only $a_A$ is committed. The same convention applies to $B$ and $C$.

\texttt{BLAKE3} compresses $64$ bytes to $32$. Each 256-bit input is split into two independently addressed 128-bit chunks: $a$ begins at $\loc{o_{A,0}}$ and $\loc{o_{A,1}}$, and $b$ at $\loc{o_{B,0}}$ and $\loc{o_{B,1}}$. Each chunk contains two consecutive $\K$ words, so its second address is the virtual multiple by $\gen$. The chaining value and digest are four consecutive words beginning at $\loc{o_{CV}}$ and $\loc{o_C}$, with their remaining addresses virtual multiples by $\gen,\gen^2,\gen^3$. Each row therefore commits six starting addresses and accesses sixteen memory words. The two metadata words encode the counter and packed block-length/flags. The BLAKE3 table (\S\ref{sec:tab-blake3}) connects them to the memory, state, and bytecode interactions, while Flock~\cite{flock} proves the compression relation.

\texttt{DEREF} stores, at the dereferenced address $\loc{\alpha}\cdot\beta$, a value chosen by a \emph{store mode} $s\in\{\texttt{deref\_cell}[\gamma],\texttt{deref\_pc},\texttt{deref\_fp}\}$:
\[
  \srcsel{s}\;=\;
  \begin{cases}
    \loc{\gamma} & s=\texttt{deref\_cell}[\gamma] \quad(\text{the local cell at offset }\gamma),\\
    \gen^{2}\cdot\pc & s=\texttt{deref\_pc} \quad(\text{a return address: the program counter advanced by two}),\\
    \fp & s=\texttt{deref\_fp} \quad(\text{the current frame pointer}).
  \end{cases}
\]
The \texttt{deref\_pc} and \texttt{deref\_fp} modes let a caller save its return address and frame so a callee can later return and restore them, the basis for function calls. The skip of two is the call's return target: place the \texttt{DEREF} that saves $\gen^{2}\cdot\pc$ immediately before the \texttt{JUMP} that transfers control, and $\gen^{2}\cdot\pc$ names the instruction just after that \texttt{JUMP}.

\texttt{DEREF\_WIDE} equates the two-word ($w=0$, internally \texttt{DEREF\_128}) or three-word ($w=1$, \texttt{DEREF\_192}) heap run beginning at $\loc{\alpha}\beta$ with the local run beginning at $\loc{\gamma}$. The width is read from public bytecode. Only the base addresses are committed; successors are virtual multiples by $\gen$ and $\gen^2$. In two-word mode the third accesses remain read-only memory-bus accesses but are excluded from the equality. Write-once witness generation fills whichever active run is unset.

\texttt{JUMP} reads a condition $c=\loc{o_c}\in\K$ and branches on whether it is zero. When $c\neq0$ it transfers control, $\pc\gets\loc{o_d}$ and $\fp\gets\loc{o_f}$; when $c=0$ it falls through, $\pc\gets\gen\cdot\pc$ and $\fp\gets\fp$.

\paragraph{Words used as addresses.} Memory words, pointers, $\pc$, and $\fp$ all share $\K$, so \texttt{DEREF} and \texttt{JUMP} require no separate tower-limb range condition.


\section{Proving primitives}\label{sec:prelim}

Every physical committed column is $\K$-valued (\S\ref{sec:vm}); every challenge, sumcheck, zerocheck, GKR value, and evaluation claim is $\E$-valued, where
\[
  \K=\F_2[x]/(x^{64}+x^4+x^3+x+1),
  \qquad
  \E=\K[y]/(y^3+y+1)\cong\F_{2^{192}}.
\]
Thus $|\E|=2^{192}$. The embedding and mixed product are
\[
  \K\hookrightarrow\E,\quad e\mapsto e+0\cdot y+0\cdot y^2,
  \qquad
  k(e_0+e_1y+e_2y^2)=ke_0+(ke_1)y+(ke_2)y^2.
\]
Committed columns, memory words, addresses, registers, and counters are $\K$-valued; their multilinear extensions are evaluated at $\E$-points and protocol randomness is sampled from $\E$. A logical extension value uses three $\K$ words/columns $V_0,V_1,V_2$ and is reassembled as $V_0+V_1y+V_2y^2$ only where an extension constraint needs it.

\textbf{Multilinear extension (MLE).} A function $b:\cube{\ell}\to\K$ (a committed column) or $b:\cube{\ell}\to\E$ (a product-tree layer of \S\ref{sec:gkr}) has a unique multilinear extension $\widetilde b:\E^{\ell}\to\E$, the polynomial of degree at most $1$ in each variable that agrees with $b$ on the boolean hypercube. Define $\eq(r,X)=\prod_i(1+r_i+X_i)$. For Boolean $r$, this is $1$ at $X=r$ and $0$ at the other Boolean points. For arbitrary $r\in\E^\ell$, it gives the interpolation weights: $\widetilde b(r)=\sum_{x\in\cube\ell}b(x)\eq(r,x)$.

\textbf{Tables and columns.} A \emph{table} is a group of columns of the same size. Each column is \emph{committed} or \emph{virtual}: a virtual column is a $\K$-linear function of the committed columns of its table, carried with no commitment of its own. Such a map, extended $\E$-linearly, commutes with the multilinear extension, $\widetilde{L(x)}(\zeta)=L\bigl(\widetilde x(\zeta)\bigr)$, so one opening of the committed columns evaluates the virtual one.

\textbf{Sumcheck.} Sumcheck reduces a claim $\sum_{x\in\cube\ell}P(x)=s$ to a single evaluation $P(\rho)$ at a random $\rho\in\E^{\ell}$, over $\ell$ rounds; a false claim passes with probability at most $\deg P\cdot\ell/|\E|$.

\textbf{Zerocheck.} To prove $P$ vanishes on $\cube\ell$, the verifier samples $r\in\E^{\ell}$ and runs sumcheck on $\sum_{x}P(x)\,\eq(r,x)=0$.

\textbf{Inner-product PCS.} An \emph{inner-product commitment scheme} lets a prover commit to a multilinear $g:\cube{M}\to\K$ and later prove a weighted sum $\sum_{w\in\cube M} W(w)\,\widehat g(w)=c$ against any weight $W:\cube{M}\to\E$ the verifier can evaluate: the commitment is over $\K$, the opening over $\E$. A point evaluation $\widehat g(\mathbf r)$ at $\mathbf r\in\E^{M}$ is the special case $W=\eq(\mathbf r,\cdot)$. WHIR and BaseFold/Ligerito are such schemes; we use one (\S\ref{sec:rs-dense}).

\section{Committing the witness}\label{sec:rs}

\subsection{Jagged commitment layout}\label{sec:stacking}

No ordinary trace column is committed on its own.  We use the Jagged PCS construction of Hemo et al.~\cite{jagged-pcs}.  Compatible equal-height columns are partitioned into power-of-two-width blocks (the ``fancy Jagged'' layout of that work): compatibility means that the columns occur in exactly the same opening groups.  For a block of width $B=2^b$ and real height $h$, its columns are stored row-major,
\[
  (P_0[0],\ldots,P_{B-1}[0],P_0[1],\ldots,P_{B-1}[h-1]).
\]
The blocks themselves are tightly concatenated and the final vector is zero-padded only once, to $2^M$.  A non-power-of-two compatible group is decomposed into power-of-two blocks, so this table packing adds no dummy cells.  Opcode columns use their announced row counts.  The memory image and memory-finalize count use the announced used-memory prefix, with public suffixes $0$ and $1$; the bytecode-finalize count uses the program-bound real instruction prefix, with public suffix $1$.  Thus $q_{\rm pkd}$ is the only full power-of-two committed column.  If $A$ is the sum of the physical block lengths, then
\[
  M=\max\!\left(15,\lceil\log_2\max(A,1)\rceil,
                    \max_{\text{blocks}}(k+b)\right),
\]
where $k$ is a block's logical row dimension and $b$ its selector width.  The last term ensures that every logical point embeds in the dense cube even when its physical prefix is short.

If such a block begins at dense offset $t$, its column $c$ and row $z$ correspond to the Boolean dense index
\[
  w=t+c+Bz \quad\text{with}\quad w<t+Bh.
\]
The low $b$ coordinates of the logical point are fixed to the bits of $c$ and the remaining coordinates are the row point.  Let $J_{t,Bh}((c,z),w)$ be the multilinear extension of this predicate.  It is the width-four read-once branching program from the Jagged PCS construction: its state records the carry of $t+(c,z)$ and the most-significant comparison seen between $w$ and $t+Bh$, while reading bits least-significant first.  The real part of a column evaluation is consequently the inner-product claim
\[
  \sum_{w\in\cube M} J_{t,Bh}((c,\zeta),w)\,\widehat q(w).
\]
The arithmetization still regards a column as extended to its power-of-two logical height with a fixed public pad value $p_i$.  The verifier adds the omitted suffix itself: if $H_i(\zeta)$ is the MLE of $[z<h_i]$, the logical claim is the real contribution plus $p_i(1-H_i(\zeta))$.  (Addition and subtraction coincide in characteristic two.)

The packed flock witness $q_{\mathrm{pkd}}$ is the sole aligned exception.  It occupies the offset-zero prefix so its ring-switch and fixed-slot claims retain their direct subcube formulas; all ordinary trace columns follow it in the tight Jagged layout.

\paragraph{Batching.} Every claim the protocol raises against a column is such a weighted sum $\sum_w W_j(w)\,\widehat q(w)=c_j$. The verifier folds the $J$ of them with powers of one random $\lambda$ into one weight $W_\lambda=\sum_j\lambda^j W_j$ and target $C_\lambda=\sum_j\lambda^j c_j$, and the inner-product PCS (\S\ref{sec:rs-dense}) proves
\[
  \sum_{w\in\cube{M}} W_\lambda(w)\,\widehat q(w)=C_\lambda.
\]
Claims covering a complete width-$B$ block receive consecutive powers in column-selector order.  Their sum is one Jagged indicator whose selector-coordinate weights are the unnormalised pairs $(1,\lambda^{2^r})$, scaled by the first power assigned to the block.  Thus the verifier runs one branching program per claimed block rather than per physical column, with no field inversions.  Unlike the generic adapter in~\cite{jagged-pcs}, our underlying PCS already accepts a verifier-evaluable inner-product weight, so we evaluate these batched indicators directly and introduce no separate Jagged reduction sumcheck.

\subsection{Inner-product PCS}\label{sec:rs-dense}

It is Ligerito~\cite{cryptoeprint:2025/1187}, instantiated by default at rate $1/2$ in the Johnson list-decoding regime with a per-level optimized slack $\eta$ and $128$-bit round-by-round soundness. A proof may instead select rate $1/2^r$ for $1\le r\le4$; the verifier checks and transcript-binds this choice. The level-zero opening claim is sampled after its commitment and therefore also binds one element of the nearby-codeword list; no separate out-of-domain evaluation is needed there. Each recursive commitment is bound by one fresh out-of-domain evaluation. Since all challenges lie in $\E=\F_{2^{192}}$, the proximity-gap and binding errors clear $128$ bits without proof-of-work grinding; the query phase grinds $17$ bits per level, so queries only close the remaining $111$. The witness is Reed--Solomon encoded over $\K$; level-zero symbols are $8$ bytes and $2^6$ interleaved symbols form each $512$-byte Merkle leaf. Openings are over $\E$: the level-zero fold pairs $\K$-symbols with $\E$-challenges, and all deeper levels are $\E$-valued.

\section{M3 model}\label{sec:m3}

The arithmetization follows the M3 model (multi-multiset matching)~\cite{m3}. The prover chooses each table's size, rows are unordered, and an instance is accepted exactly when every table's constraints hold and the bus balances.

\paragraph{Constraints.} A constraint is a polynomial in the columns that must vanish on every row, including the padding rows.

\paragraph{The bus.} All interactions share a single \emph{bus}, a channel carrying tuples of $m$ field elements. A table flushes by wiring its columns into a tuple and either \emph{pushing} it onto the bus or \emph{pulling} it off, once per row. A few fixed \emph{boundary} tuples, belonging to no table, are pushed or pulled directly. The width $m$ is that of the widest tuple, twelve (the bytecode tuple: separator, $\pc$, count, opcode, and eight operand/immediate slots); shorter tuples are zero-padded.

\paragraph{Domain separation.} The first coordinate of every tuple is a \emph{domain separator}, a fixed constant naming its interaction: $\dsep{ST}$ for state, $\dsep{MEM}$ for memory, $\dsep{BC}$ for bytecode. Since everything shares the one bus, the separators are what keep the three kinds of tuple from cancelling against one another: a memory tuple can never meet a bytecode tuple. The interactions below are three uses of this one channel, told apart by this coordinate.

\paragraph{Balance.} The bus balances when its pushed tuples and its pulled tuples form the same multiset: the same tuples with the same multiplicities, in any order. This is the instance's only global check, proven by grand product argument (\S\ref{sec:gp}).

\paragraph{Per-instruction tables.} Each instruction has its own table, with columns $\pc$, $\fp$, $\mathit{next\_pc}$, $\mathit{next\_fp}$, the operands, the operand addresses, and the read values. For every instruction except \texttt{JUMP}, the successor is virtual: $\mathit{next\_fp}=\fp$ is a copy of $\fp$, and $\mathit{next\_pc}=\gen\cdot\pc$ is the free increment of \S\ref{sec:vm}. Each row makes the same kinds of flush: on the state interaction (\S\ref{sec:state}) it pulls $\tup{\dsep{ST},\pc,\fp}$ and pushes $\tup{\dsep{ST},\mathit{next\_pc},\mathit{next\_fp}}$; it makes one memory flush per cell read (\S\ref{sec:memchan}) and one bytecode flush for the instruction at $\pc$ (\S\ref{sec:bytecode}). The concrete tables are in \S\ref{sec:tables}.

\subsection{Proving the local constraints}\label{sec:air}

Each table must satisfy its own constraints on every one of its own rows. Instead of one zerocheck per table, over cubes of different sizes, we batch them into a single sumcheck.

Table $T_j$ has $2^{\tau_j}$ rows and constraints $C_{j,1},\dots,C_{j,s_j}$: polynomials of degree at most $d$ in a row's column values, required to vanish on every row ($d=2$ throughout). Write $\sigma=\sum_j s_j$ for the total number of constraints and $n=\max_j\tau_j$ for the tallest table's height.

\paragraph{One challenge and one point.} The verifier samples $\eta\in\E$ and fixes a point $r\in\E^{n}$ (in the assembled protocol $r$ is not freshly sampled but taken from the bus, see \emph{Nonzero claims} below and \S\ref{sec:leafstack}). Table $T_j$ mixes its constraints with the powers of $\eta$ from $o_j=s_1+\dots+s_{j-1}$ on, and is tested at the length-$\tau_j$ prefix of $r$:
\[
  C_j\;=\;\sum_{i=1}^{s_j}\eta^{\,o_j+i-1}\,C_{j,i},
  \qquad
  S_j\;=\;\widetilde{C_j}\bigl(r_{<\tau_j}\bigr)\;=\;\sum_{x\in\cube{\tau_j}}C_j(x)\,\eq(r_{<\tau_j},x).
\]
The ranges of $\eta$-powers are \emph{disjoint}, so no table's violation can be cancelled by another's, and proving $\sum_j S_j=0$ proves every $C_{j,i}$ vanishes on every row of its own table, up to Schwartz--Zippel error $(n+\sigma)/|\E|$ over $r$ and $\eta$. (The identities keep disjoint ranges throughout. \S\ref{sec:leafstack} appends three more entries per table that deliberately \emph{share} their powers, since only their total over tables has to be pinned; that lengthens the $\eta$-degree, and the error, by $3$.)

\paragraph{One sumcheck.} The $S_j$ are sumcheck claims of different lengths. Lift each onto the common cube $\cube n$ by attaching the indicator of the all-ones assignment on the $n-\tau_j$ variables $T_j$ does not use:
\[
  f_j(X_0,\dots,X_{n-1})\;=\;\bigl(C_j\,\eq(r_{<\tau_j},\cdot)\bigr)(X_0,\dots,X_{\tau_j-1})\cdot\prod_{i\ge\tau_j}X_i .
\]
Because $\sum_x\prod_i x_i=1$ on any cube, $\sum_{x\in\cube n}f_j(x)=S_j$, so with $F=\sum_j f_j$ the whole statement is the single claim $\sum_{x\in\cube n}F(x)=0$, proven by $n$ rounds of sumcheck on $F$.

The rounds bind $X_{n-1}$ first and $X_0$ last, so the first $n-\tau_j$ touch only variables $T_j$ does not use: there $f_j$ contributes the line $Y\,k\,S_j$, with $k$ the product of the challenges drawn so far (the same for every table that is waiting) and $S_j=0$. Table $T_j$ enters at round $n-\tau_j$ and binds $X_{\tau_j-1},\dots,X_0$ in turn. This is back-loaded batching: a short claim waits, then joins weighted by the rounds it sat out.

\paragraph{Nonzero claims.} Nothing above needs $S_j=0$, nor $r$ to be a fresh challenge. Any polynomial in $T_j$'s columns whose $\eq$-weighted sum over the rows is a value the verifier knows is batched exactly like a constraint, becoming one more entry of $C_j$ with a power of $\eta$ of its own; only $S_j$ changes, from $0$ to that value. And any point fixed at random after the trace commitment serves as $r$. \S\ref{sec:leafstack} does both, running the batch at $r=\zeta$, the point GKR leaves behind, and handing each table three such polynomials. What is proven is then
\[
  \sum_{x\in\cube n}F(x)\;=\;\sum_j S_j .
\]

The one consequence for the rounds is that a waiting table's line $Y\,k\,S_j$ no longer vanishes. Write $X_\ell$ for the variable a round binds, $Y$ for it while still an unknown, and $k$ for the product of the challenges drawn so far, which is common to every waiting table. They therefore sum to $Y\,u$ for the single scalar $u=k\sum_{\tau_j\le\ell}S_j$, and with $p$ the joined tables' cofactor, of degree at most $2$, the round polynomial is the cubic
\[
  h(Y)\;=\;\eq(r_\ell,Y)\,p(Y)+Y\,u .
\]
The prover sends $h$ whole, at $\{0,1,\gen,\gen^{2}\}$: four field elements a round, one more than \S\ref{sec:gkr}'s $\eq$-trick. That one is the price of the nonzero claims rather than of sending $h$ whole, for the trick cannot save it: the verifier has no way to form $u$, so $u$ would travel beside the cofactor instead. The rounds are then vanilla sumcheck, checking $h(0)+h(1)=c$ against the running claim $c$ and folding $c\leftarrow h(\rho_\ell)$, with no $\eq$ reapplied to the message. Prover work is unchanged: $p$ is quadratic, so its fourth value interpolates from the other three rather than costing another pass over the rows.

\subsection{Balancing the bus: grand product}\label{sec:gp}

The bus balances when its pushed and pulled tuples form the same multiset (\S\ref{sec:m3}). Two random challenges, $\alpha$ then $\gamma$, reduce this to a single product.

\emph{Fingerprint.} The verifier samples $\alpha\in\E$ and maps each width-$m$ tuple $\sigma=(\sigma_0,\dots,\sigma_{m-1})$ to a single $\E$-element,
\[
\pi_\alpha(\sigma)\;=\;\sum_{i=0}^{m-1}\alpha^{i}\,\sigma_i ,
\]
each coordinate $\sigma_i$ a committed column value, a virtual column value, or a public constant, all $\K$-elements, so each summand is a mixed $\K\times\E$ product.

\emph{Product check.} The verifier samples $\gamma\in\E$ and the prover proves
\begin{equation}
\prod_{\sigma\ \text{pushed}}\bigl(\gamma-\pi_\alpha(\sigma)\bigr)
\;=\;
\prod_{\sigma\ \text{pulled}}\bigl(\gamma-\pi_\alpha(\sigma)\bigr).
\label{eq:gp}
\end{equation}
\emph{Completeness and soundness.} If the pushed and pulled multisets agree, the two products in \eqref{eq:gp} have the same factors and are equal for every $\alpha,\gamma$; completeness is immediate. For soundness, read each side as the polynomial $\prod_\sigma\bigl(X-\pi_\alpha(\sigma)\bigr)$ evaluated at $X=\gamma$: its roots are that side's fingerprints, which over the formal $\alpha$ are the tuples themselves, so the polynomial determines that side's multiset. The two sides therefore agree as polynomials in $(\alpha,\gamma)$ exactly when the multisets do. When they differ, Schwartz--Zippel applies jointly to the tuple-fingerprinting challenge $\alpha$, the product challenge $\gamma$, the padding correction, and the GKR reductions. If each side has at most $2^\mu$ factors and tuples have width $m$, a conservative total degree bound is
\[
  2m\,2^\mu+8(\mu+1)^2.
\]
The verifier checks that this bound over $|\E|=2^{192}$ leaves at least $128$ bits of soundness. The count product is verified by the same GKR argument, and its root is checked directly for nonzero; it introduces no separate root-at-random test. No Fiat--Shamir grinding is needed for the grand-product phase.

Each side's product is proven by a GKR pass over a product tree (\S\ref{sec:gkr}).

\paragraph{Padding.} Every table is padded past its real rows to a power of two. Each padding row is a fixed default: every column zero but the read counts, set to $\gen^{0}=1$. Its memory and bytecode flushes are then default tuples that do not self-cancel, so the prover announces each table's real row count and the verifier divides the surplus out of \eqref{eq:gp}: a default tuple $d$ in surplus $\delta_d$ on a side divides that side by $(\gamma-\pi_\alpha(d))^{\delta_d}$.


\subsection{Proving the grand product (GKR)}\label{sec:gkr}

Take one side of \eqref{eq:gp}. Its tuples give leaves $v_k=\gamma-\pi_\alpha(\sigma_k)$, padded with $1$s to a power of two, $2^\mu$ in all. Their product $P=\prod_kv_k$ is the root of a binary product tree: layer $0$ is the leaves, layer $i$ has $2^{\mu-i}$ nodes, and layer $\mu$ is $P$. Let $\widetilde V_i:\E^{\mu-i}\to\E$ denote the multilinear extension of layer $i$.

GKR contracts two binary levels at a time. Four grandchildren satisfy
\[
V_i(x)=\prod_{a,b\in\{0,1\}}V_{i-2}(a,b,x),
\]
and hence
\[
\widetilde V_i(r)=\sum_{x\in\cube{\mu-i}}\eq(r,x)\prod_{a,b\in\{0,1\}}\widetilde V_{i-2}(a,b,x).
\]
A sumcheck reduces a claim $\widetilde V_i(r)$ to the four values $\widetilde V_{i-2}(a,b,\rho)$ at one random point $\rho$. The product part of each round polynomial has degree four. As in the usual $\eq$-trick~\cite{cryptoeprint:2024/108}, the known equality factor is not transmitted. Moreover, the incoming claim already fixes one coefficient. Write the degree-four cofactor as $q(X)=\sum_{j=0}^{4}c_jX^j$, let $s$ be the current coordinate of $r$, and let $T$ be the incoming claim. The prover sends only
\[
d=q(0)+q(1),\qquad c_2,\qquad c_3,\qquad c_4.
\]
Since $T=(1+s)q(0)+sq(1)=c_0+sd$, the verifier recovers $c_0=T+sd$ and $c_1=d+c_2+c_3+c_4$, then evaluates $q$ at the new challenge by Horner's rule. Thus each round sends exactly the four independent field elements permitted by its consistency equation.

At the end of the layer the prover sends the four grandchildren. Two fresh challenges $(u_0,u_1)$ combine them multilinearly, leaving one claim $\widetilde V_{i-2}(u_0,u_1,\rho)$. If $\mu$ is odd, GKR first processes the root-most binary layer, which has no sumcheck rounds, and then proceeds entirely in radix four. After $\lceil\mu/2\rceil$ layers a single leaf claim $\widetilde V_0(\zeta)$ remains, which \S\ref{sec:leafstack} settles against the committed columns. Prover work remains $O(2^\mu)$ and the proof remains $O(\mu^2)$ field elements, with fewer layers and a smaller constant.

\paragraph{Batching the three trees.} The bus needs three grand products: push, pull, and counts. They are batched in a single GKR pass (verifier-friendliness, good for recursion). Push and pull have equal depth $\mu$ (their blocks come in matched pairs); the count tree, of depth $\mu_c\le\mu$, is padded with $1$-leaves to $2^{\mu}$.

The verifier samples a combiner $\lambda$, and each radix-four layer runs one sumcheck on the combined identity
\[
\sum_t\lambda^t\widetilde V_i^t(r)=\sum_{x\in\cube{\mu-i}}\eq(r,x)\sum_t\lambda^t\prod_{a,b\in\{0,1\}}\widetilde V_{i-2}^t(a,b,x),\qquad t\in\{0,1,2\}.
\]
A fresh $\lambda$ after every layer binds the three individual tail values rather than only their previous linear combination. Consequently push, pull, and count all reduce to claims at the same point $\zeta$.

The identity padding is logical: the prover stores only the count tree's actual power-of-two prefix. Once that prefix has folded to one value, each remaining gate pairs it with $1$, while every omitted all-one gate is known and contributes nothing to a nonconstant sumcheck message. The prover also stores each radix-four layer in its original interleaved child order and constructs only the even-numbered binary levels consumed by the protocol. Thus the radix-four verifier savings do not require materializing the padded tree or transposing four child tables.

\subsection{Stacking the leaves}\label{sec:leafstack}

GKR (\S\ref{sec:gkr}) reduces one side of \eqref{eq:gp} to a single evaluation $\widetilde V_0(\zeta)$ of its \emph{leaves}, the product's factors $\gamma-\pi_\alpha(\sigma)$, one per bus tuple $\sigma$. We settle this evaluation against the committed columns.

\paragraph{Flush blocks.} The leaves split into \emph{blocks}, one per flush rule and one per boundary set. A block of a $2^{\kappa}$-row table holds $2^{\kappa}$ leaves; the row-$z$ leaf is $\gamma-\sum_{i=0}^{m-1}\alpha^{i}\,c_{i}(z)$, with $c_{i}(z)$ coordinate $i$ of the row's tuple. Each $c_i$ is public, a committed column, or a virtual column.

\paragraph{Layout.} Order the blocks largest first at aligned offsets, padding to $2^{\mu}$ leaves with $1$s, which leave the product unchanged.  This is deliberately the older aligned block layout, not the Jagged commitment layout of \S\ref{sec:stacking}: Jagged is used only for the committed trace columns.  Block $b$, of $2^{\kappa_b}$ leaves, occupies the subcube $\{(z,\mathsf{sel}_b):z\in\cube{\kappa_b}\}$ for a fixed selector $\mathsf{sel}_b$.  The grand-product arguments and their polynomial stacking are otherwise unchanged.

\paragraph{Decomposition.} Split $\zeta=(\zeta^{b}_{\mathrm{lo}},\zeta^{b}_{\mathrm{hi}})$ at block $b$'s boundary. Each block is a subcube, so the leaf MLE splits into a public selector times the block's own MLE, the padding contributing the leftover mass:
\[
  \widetilde V_0(\zeta)\;=\;\sum_b \eq(\mathsf{sel}_b,\zeta^{b}_{\mathrm{hi}})\Bigl(\gamma-\sum_{i=0}^{m-1}\alpha^{i}\,\widetilde c_{b,i}(\zeta^{b}_{\mathrm{lo}})\Bigr)\;+\;\Bigl(1-\sum_b\eq(\mathsf{sel}_b,\zeta^{b}_{\mathrm{hi}})\Bigr).
\]
The verifier forms the selectors $\eq(\mathsf{sel}_b,\cdot)$ and the constant coordinates itself. Each remaining $\widetilde c_{b,i}(\zeta^{b}_{\mathrm{lo}})$ is settled by its kind: a committed column by a direct opening (\S\ref{sec:prelim}), a virtual column off the committed columns it images (so $\gen\cdot\pc$ adds nothing beyond the $\pc$ opening), a \texttt{BLAKE3} value coordinate by a point evaluation of $q_\pkd$ at its slot (\S\ref{sec:tab-blake3}), the index column by \S\ref{sec:idxcol}. A committed column feeding several equal-size blocks is opened once, at their shared point $\zeta_{\mathrm{lo}}$.

\paragraph{Blocks of a table.} Split that sum by owner. Three blocks per side belong to no table, the state boundary and the memory and bytecode seeds (finalizations on the pull side); they settle exactly as above, and the counts side has none of them. Every other block belongs to a table, and a table owns many: one per flush on push and on pull, one per read count on the counts, so \texttt{BLAKE3} owns ten, ten and nine.

All of $T_j$'s blocks are $2^{\tau_j}$ rows tall, so they read its columns at the one point $\zeta_{<\tau_j}$, and that shared point is what collapses them. Fix a side, let $\mathcal B$ be the blocks $T_j$ owns there, and put the row values back where those evaluations stand:
\[
  \sum_{b\in\mathcal B}\eq(\mathsf{sel}_b,\zeta^{b}_{\mathrm{hi}})\Bigl(\gamma-\sum_{i=0}^{m-1}\alpha^{i}\,\widetilde c_{b,i}(\zeta_{<\tau_j})\Bigr)
  \;=\;\sum_{x\in\cube{\tau_j}}\eq(\zeta_{<\tau_j},x)\,B_j(x),
\]
where $B_j$ is that same expression read at a row $x$ of $T_j$, a polynomial every coefficient of which the verifier forms itself (on the counts side $\alpha=1$ and $\gamma=0$, \S\ref{sec:memchan}, so it is a bare sum of count columns). The right-hand side is exactly a claim of \S\ref{sec:air}'s kind, with a nonzero value. So hand the three $B_j$, one per side, to the batch and run it at $r=\zeta$. They take three powers of $\eta$ placed after all $\sigma=\sum_j s_j$ identity powers of \S\ref{sec:air} and \emph{shared} by every table, so with a table's real constraints vanishing,
\[
  S_j\;=\;\eta^{\,\sigma}\,\widetilde B^{\mathrm{push}}_j+\eta^{\,\sigma+1}\,\widetilde B^{\mathrm{pull}}_j+\eta^{\,\sigma+2}\,\widetilde B^{\mathrm{cnt}}_j ,
\]
all at $\zeta_{<\tau_j}$.

Sharing those powers is what ties the batch to the bus. Summing over the tables,
\[
  \sum_j S_j\;=\;\eta^{\,\sigma}R_{\mathrm{push}}+\eta^{\,\sigma+1}R_{\mathrm{pull}}+\eta^{\,\sigma+2}R_{\mathrm{cnt}} ,
\]
each $R_{\mathrm{side}}$ being that side's leaf claim less its table-less blocks and padding, which the verifier forms for itself. The batch therefore starts from a target it derived, and $\eta$ is drawn only afterwards, so reaching that one number forces the three bus equations $\sum_j\widetilde B^{\mathrm{side}}_j=R_{\mathrm{side}}$ at once, and zero in every identity slot besides. Per-table powers would not factor this way: the verifier would need three shares per table transmitted, which is sound but costly, and a transmitted aggregate would settle nothing at all, appearing in no other check.

So the bus sends nothing for a table's blocks, and no column of a table is opened at $\zeta$: no such evaluation travels at all. The batch leaves one claim per column of $T_j$ instead, at its own output point $\rho_{<\tau_j}$, settled as above.

\subsection{State interaction}\label{sec:state}

The state interaction carries $\tup{\dsep{ST},\pc,\fp}$. Each instruction row pulls its current state and pushes its successor $\tup{\dsep{ST},\mathit{next\_pc},\mathit{next\_fp}}$; the bus is preseeded by pushing the initial state $\tup{\dsep{ST},\pc_0,\fp_0}$ and pulling the final state $\tup{\dsep{ST},\pc_{\mathrm{final}},\fp_{\mathrm{final}}}$. When the bus balances, the per-row transitions form a chain from $(\pc_0,\fp_0)$ to $(\pc_{\mathrm{final}},\fp_{\mathrm{final}})$, possibly alongside disjoint cycles. The chain is the execution; a stray cycle is a closed loop of rows that still satisfies every constraint and reads the same committed memory and bytecode, so it cannot affect the start-to-final claim. By convention execution halts at a fixed sentinel: the last bytecode slot. Writing $N$ for the program length (a power of two, \S\ref{sec:e2e-bc}), the instructions sit at $\gen^{0},\dots,\gen^{\,N-1}$, so $\pc_{\mathrm{final}}=\gen^{\,N-1}$, with $\fp_{\mathrm{final}}=\gen^{0}$. A compiled \texttt{main} ends by jumping there. Both boundary states are then public, so the verifier forms them itself.

\section{Memory and bytecode lookups}\label{sec:omc}

Both a memory read and a bytecode read are indexed lookups into an array: the committed memory, or the public program. We realize each as an interaction on the shared bus (\S\ref{sec:m3}), balanced by \eqref{eq:gp}, using offline memory checking with a read count per index.

\subsection{Memory interaction}\label{sec:memchan}

The interaction carries $\tup{\dsep{MEM},\addr,\cnt,\val}$, with every
coordinate in $\K$. The prover commits one memory-value column $M$ and the
counts (powers of $\gen$). With $A[i]$ the number of reads of address $i$, the
flushes are:
\begin{itemize}
\item \textbf{Initialize:} for each address $i$, \push\ $\tup{\dsep{MEM},\gen^{i},1,\mem[i]}$ (seed count $\gen^{0}=1$).
\item \textbf{Read} (value $v$ at address $a$, count $\cnt=\gen^{t}$): \pull\ $\tup{\dsep{MEM},a,\cnt,v}$ and \push\ $\tup{\dsep{MEM},a,\gen\cdot\cnt,v}$.
\item \textbf{Finalize:} for each address $i$, \pull\ $\tup{\dsep{MEM},\gen^{i},\gen^{A[i]},\mem[i]}$, the committed final count $\gen^{A[i]}$.
\end{itemize}
\paragraph{Counts must be nonzero.} A read multiplies its count by $\gen$, pulling $\tup{\dsep{MEM},a,\cnt,v}$ and pushing $\tup{\dsep{MEM},a,\gen\cdot\cnt,v}$. If $\cnt=0$ these coincide (as $\gen\cdot0=0$) and cancel, so the read disappears from the bus and its $v$ is never checked against $\mem[a]$. Every count must therefore be nonzero, i.e.\ a power of $\gen$. A third grand product, over all count columns (memory and bytecode), enforces this: its root is nonzero iff every count is (\S\ref{sec:gkr}).

\paragraph{Counts must not wrap.} A count is the power $\gen^{t}$ after $t$ reads, so the total number of reads must stay below $\operatorname{ord}(\gen)=2^{64}-1$. At $64$ bits this is not waved through as automatic. The protocol fixes public \emph{instance caps}, and the verifier checks the announced instance against them before running any reduction (\S\ref{sec:e2e-unrolled}): memory size $16\le h\le 32$ (\S\ref{sec:vm}), each table's row count below $2^{32}$, and bytecode length at most $2^{32}$. A row makes at most $17$ reads (the \texttt{BLAKE3} row: sixteen memory, one bytecode), so every \emph{admissible} instance performs
\[
  R\;\le\;17\cdot 9\cdot 2^{32}\;<\;2^{40}
\]
read flushes in total, far below $2^{64}-1$: the powers $\gen^{0},\gen^{1},\dots$ occurring in one address's history are pairwise distinct, and the counting arguments of this section are exact, with no probabilistic term. The caps cost no completeness ($2^{32}$ rows per table is far beyond any provable trace).

\paragraph{What balance certifies.} Fix an address $i$. Its only non-read flushes are the seed, a push of count $\gen^{0}$ with value $\mem[i]$, and the finalize, a pull of the committed final count with value $\mem[i]$.

\emph{Completeness.} An honest prover gives the $A[i]$ reads of $i$ the counts $\gen^{0},\dots,\gen^{A[i]-1}$ and the finalize the count $\gen^{A[i]}$. The read at count $\gen^{t}$ pulls $\gen^{t}$ and pushes $\gen^{t+1}$, so both sides are the multiset $\{\gen^{0},\dots,\gen^{A[i]}\}$ at value $\mem[i]$, and the bus balances.

\emph{Soundness.} Conversely, suppose the bus balances; we show every read of $i$ returns $\mem[i]$. Take any value $v\neq\mem[i]$. The seed and finalize carry $\mem[i]$, so the value-$v$ tuples come only from reads returning $v$, each pulling its count $\cnt$ and pushing $\gen\cdot\cnt$. Restricted to value $v$, balance equates the pulled with the pushed counts:
\[
  \{\,\cnt : \text{reads returning } v\,\}\;=\;\{\,\gen\cdot\cnt : \text{reads returning } v\,\}.
\]
This multiset of counts is thus invariant under multiplication by $\gen$. But the counts are powers of $\gen$, and $\times\gen$ cycles through all $2^{64}-1$ of them, so an invariant multiset weights every power equally: its size is a multiple of $2^{64}-1$. Under the instance caps its size is at most $R<2^{39}$, hence it is empty. So no read returns $v$, and every read of $i$ returns $\mem[i]$.

\subsection{Bytecode interaction}\label{sec:bytecode}

Instruction fetch is the same as the memory construction (\S\ref{sec:memchan}) with the \emph{public} program for the committed memory and the program counter for the address. The tuple is $\tup{\dsep{BC},\pc,\cnt,\mathit{instr}}$, where $\mathit{instr}$ is the encoded instruction: an opcode and eight operand/immediate slots. Shorter instructions zero their unused high slots.

\subsection{The index column}\label{sec:idxcol}

The memory seed and finalize tuples (\S\ref{sec:memchan}) run over every address $\gen^{i}$, $i=0,\dots,2^{h}-1$. Settling the grand product leaves the verifier evaluating, at a random $\zeta\in\E^{h}$, the MLE of this \emph{index column} $\gen^{0},\gen^{1},\dots,\gen^{2^{h}-1}$. It is not committed: as $\gen$-powers the MLE factors, costing $O(h)$.

Place address $\gen^{i}$ at the cube point $w\in\cube h$ of $i$'s bits, so $i=\sum_k 2^{k}w_k$ and
\[
  \gen^{\,i}\;=\;\prod_{k=0}^{h-1}\bigl(\gen^{\,2^{k}}\bigr)^{w_k}.
\]
Its MLE is the same product with each bit $w_k$ relaxed to $\zeta_k$,
\[
  \widetilde{\mathrm{idx}}(\zeta)\;=\;\prod_{k=0}^{h-1}\bigl(1+\zeta_k\,(1+\gen^{\,2^{k}})\bigr),
\]
each factor $1$ at $\zeta_k=0$ and $\gen^{2^{k}}$ at $\zeta_k=1$. The verifier evaluates it in $O(h)$ from the squares $\gen^{2^{0}},\gen^{2^{1}},\dots$. (Bytecode is the same, with the program counter for the address.)

\section{The instruction tables}\label{sec:tables}

Each instruction has its own table.

\subsection{\texttt{XOR}}\label{sec:tab-xor}

Asserts $\loc{o_C}=\loc{o_A}+\loc{o_B}$ in $\K$ (bitwise \textsc{xor}).
\begin{itemize}
\item \textbf{Columns:} $\pc,\fp$; operands $o_A,o_B,o_C$; addresses $a_A,a_B,a_C$; values $v_A,v_B,v_C\in\K$; memory counts $r_A,r_B,r_C$; bytecode count $r_{\mathrm{bc}}$.
\item \textbf{Constraints:} the three addresses $a_X=\fp\cdot o_X$ and the base-field sum $v_C=v_A+v_B$.
\item \textbf{Flushes:}
  \begin{itemize}
  \item \textbf{State:} \pull\ $\tup{\dsep{ST},\pc,\fp}$, \push\ $\tup{\dsep{ST},\gen\cdot\pc,\fp}$.
  \item \textbf{Bytecode:} \pull\ $\tup{\dsep{BC},\pc,r_{\mathrm{bc}},\opc{XOR},o_A,o_B,o_C,0,0,0,0,0}$, \push\ the same at $\gen\cdot r_{\mathrm{bc}}$.
  \item \textbf{Memory} ($X\in\{A,B,C\}$): \pull\ $\tup{\dsep{MEM},a_X,r_X,v_X}$, \push\ $\tup{\dsep{MEM},a_X,\gen\cdot r_X,v_X}$.
  \end{itemize}
\end{itemize}

\subsection{\texttt{MUL}}\label{sec:tab-mul}

Identical to \texttt{XOR} but for the base-field product:
\begin{itemize}
\item \textbf{Columns:} as \texttt{XOR}.
\item \textbf{Constraints:} the three addresses and $v_C=v_Av_B$ in $\K$.
\item \textbf{Flushes:} as \texttt{XOR}, with opcode $\opc{MUL}$.
\end{itemize}

\subsection{\texttt{XOR\_192} and \texttt{MUL\_192}}\label{sec:tab-ext}

Each operand is three consecutive base words. The table commits one address
$a_X=\fp o_X$ per operand and uses $\gen a_X,\gen^2a_X$ virtually for the
remaining coordinates. It performs nine memory accesses, reassembles
$v_X=v_{X,0}+v_{X,1}y+v_{X,2}y^2$, and constrains respectively
$v_C=v_A+v_B$ or $v_C=v_Av_B$ in $\E$. The multiplication table also has a
public-bytecode mode $m=1$ for $v_A\in\K$: only $v_{A,0}$ participates in the
product, while the other two memory cells remain ordinary read-only memory-bus
accesses. Two quadratic constraints tie the effective upper limbs to the
memory values when $m=0$ and to zero when $m=1$, so the AIR remains degree two.
In the zkDSL, \texttt{xor\_192} exposes addition,
\texttt{mul\_192\_base} exposes the scalar mode, and \texttt{div\_192} is a
checked multiplication with the quotient back-solved by witness generation.

\subsection{\texttt{SET\_CONSTANT}}\label{sec:tab-set}

Asserts $\loc{o}=k\in\K$, how a constant enters the read-only memory.
\begin{itemize}
\item \textbf{Columns:} $\pc,\fp$; operand $o$; immediate $k\in\K$; address $a$; memory count $r$; bytecode count $r_{\mathrm{bc}}$.
\item \textbf{Constraints:} the address $a=\fp\cdot o$.
\item \textbf{Flushes:}
  \begin{itemize}
  \item \textbf{State:} \pull\ $\tup{\dsep{ST},\pc,\fp}$, \push\ $\tup{\dsep{ST},\gen\cdot\pc,\fp}$.
  \item \textbf{Bytecode:} \pull\ $\tup{\dsep{BC},\pc,r_{\mathrm{bc}},\opc{SET},o,k,0,0,0,0,0,0}$, \push\ the same at $\gen\cdot r_{\mathrm{bc}}$.
  \item \textbf{Memory:} \pull\ $\tup{\dsep{MEM},a,r,k}$, \push\ $\tup{\dsep{MEM},a,\gen\cdot r,k}$.
  \end{itemize}
\end{itemize}

\subsection{\texttt{DEREF}}\label{sec:tab-deref}

\texttt{DEREF} follows a pointer and asserts the cell it points to. The row reads three cells:
\begin{itemize}
\item the \emph{pointer} $p$, at $a_1=\fp\cdot\alpha$;
\item the \emph{target} $v_2$, at the dereferenced address $a_2=p\cdot\beta$;
\item the \emph{local} value $v_3$, at $a_3=\fp\cdot\gamma$.
\end{itemize}
The store asserts that $v_2$ equals the source chosen by the mode $s$ (\S\ref{sec:vm}): the local cell $v_3$, the return address $\gen^{2}\cdot\pc$ (a virtual $\times\gen^{2}$ of $\pc$, \S\ref{sec:prelim}), or the frame pointer $\fp$. The mode is two boolean flags $(f_{\pc},f_{\fp})$: $(0,0)$ for $\texttt{deref\_cell}[\gamma]$, $(1,0)$ for $\texttt{deref\_pc}$, $(0,1)$ for $\texttt{deref\_fp}$.
\begin{itemize}
\item \textbf{Columns:} $\pc,\fp$; operands $\alpha,\beta,\gamma$; flags $f_{\pc},f_{\fp}$; addresses $a_1,a_2,a_3$; words $p,v_2,v_3\in\K$; memory counts $r_1,r_2,r_3$; bytecode count $r_{\mathrm{bc}}$.
\item \textbf{Constraints:} the addresses $a_1=\fp\cdot\alpha$, $a_2=p\cdot\beta$, $a_3=\fp\cdot\gamma$, and the store
\[
  v_2 \;=\; (1+f_{\pc}+f_{\fp})\,v_3 \;+\; f_{\pc}\,(\gen^{2}\cdot\pc) \;+\; f_{\fp}\,\fp ,
\]
which gives $v_3$, $\gen^{2}\cdot\pc$, or $\fp$ at the three flag settings. All relations are in $\K$. The flags need no validity check: the verifier reads them from the public program (\S\ref{sec:e2e-bc}).
\item \textbf{Flushes:}
  \begin{itemize}
  \item \textbf{State:} \pull\ $\tup{\dsep{ST},\pc,\fp}$, \push\ $\tup{\dsep{ST},\gen\cdot\pc,\fp}$.
  \item \textbf{Bytecode:} \pull\ $\tup{\dsep{BC},\pc,r_{\mathrm{bc}},\opc{DRF},\alpha,\beta,\gamma,f_{\pc},f_{\fp},0,0,0}$, \push\ the same at $\gen\cdot r_{\mathrm{bc}}$.
  \item \textbf{Memory:} three reads, each \pull\ at count $r$ and \push\ at $\gen\cdot r$: pointer $\tup{\dsep{MEM},a_1,r_1,p}$, target $\tup{\dsep{MEM},a_2,r_2,v_2}$, local cell $\tup{\dsep{MEM},a_3,r_3,v_3}$.
  \end{itemize}
\end{itemize}

\subsection{\texttt{DEREF\_WIDE}}\label{sec:tab-deref-ext}

\texttt{DEREF\_WIDE} bundles consecutive scalar cell-mode dereferences. It
reads the pointer $p$ at $a_1=\fp\alpha$, the heap run beginning at
$a_2=p\beta$, and the local run beginning at $a_3=\fp\gamma$. The successors
$\gen a_2,\gen^2a_2$ and $\gen a_3,\gen^2a_3$ are virtual columns. A width bit
$w$, fixed by public bytecode, selects native two-word \texttt{DEREF\_128}
($w=0$) or three-word \texttt{DEREF\_192} ($w=1$). Its constraints are the
three base-address relations and the packed equality
\[
  v_{2,0}+yv_{2,1}+wy^2v_{2,2}
  =v_{3,0}+yv_{3,1}+wy^2v_{3,2}.
\]
The row always flushes one pointer access and all six run-cell accesses to
memory, plus the usual state and bytecode interactions. In two-word mode the
third cells are independently memory-bound padding reads, but are not equated.
Thus one VM cycle transfers either native 128-bit data or one three-word
extension value without weakening the 64-bit memory-word invariant.

\subsection{\texttt{JUMP}}\label{sec:tab-jump}

A conditional jump on whether $c=\loc{o_c}\in\K$ is nonzero. It transfers to $\pc\gets d$, $\fp\gets f$ when $c\neq0$ and falls through otherwise. A committed boolean indicator $b=[\,c\neq0\,]$ and inverse $w\in\K$ certify
\[
  b=c\cdot w,\qquad c\,(b+1)=0 .
\]
When $c\neq0$ the second forces $b=1$ and the first then $w=c^{-1}$; when $c=0$ the first forces $b=0$. Writing $s=\gen\cdot\pc$ for the virtual fall-through, the committed next state is the selection
\[
  \mathit{next\_pc}=b\cdot d+(b+1)\,s,\qquad \mathit{next\_fp}=b\cdot f+(b+1)\,\fp.
\]
Alone among the tables, \texttt{JUMP} commits its successors.
\begin{itemize}
\item \textbf{Columns:} $\pc,\fp,\mathit{next\_pc},\mathit{next\_fp}$; operands $o_c,o_d,o_f$; addresses $a_c,a_d,a_f$; words $c,d,f,w\in\K$; indicator $b\in\{0,1\}$; memory counts $r_c,r_d,r_f$; bytecode count $r_{\mathrm{bc}}$.
\item \textbf{Constraints:} the three addresses, the indicator pair $b=cw$ and $c(b+1)=0$, and the two next-state selections above, all over $\K$.
\item \textbf{Flushes:}
  \begin{itemize}
  \item \textbf{State:} \pull\ $\tup{\dsep{ST},\pc,\fp}$, \push\ $\tup{\dsep{ST},\mathit{next\_pc},\mathit{next\_fp}}$.
  \item \textbf{Bytecode:} \pull\ $\tup{\dsep{BC},\pc,r_{\mathrm{bc}},\opc{JMP},o_c,o_d,o_f,0,0,0,0,0}$, \push\ the same at $\gen\cdot r_{\mathrm{bc}}$.
  \item \textbf{Memory} ($X\in\{c,d,f\}$): \pull\ $\tup{\dsep{MEM},a_X,r_X,X}$, \push\ $\tup{\dsep{MEM},a_X,\gen\cdot r_X,X}$.
  \end{itemize}
\end{itemize}

\subsection{\texttt{BLAKE3}}\label{sec:tab-blake3}

Compresses two 256-bit operands to a 256-bit digest. Each input is two
independently addressed 128-bit chunks. Their committed starting addresses are
$a_{A,i}=\fp o_{A,i}$ and $a_{B,i}=\fp o_{B,i}$ for $i\in\{0,1\}$; the second
word of each chunk is at the virtual address $\gen a_{X,i}$. The four-word
chaining value and output start at the committed addresses $a_{CV}=\fp o_{CV}$
and $a_C=\fp o_C$, with successor addresses $\gen^j a_{CV}$ and $\gen^j a_C$
for $j\in\{1,2,3\}$. Thus each row accesses sixteen memory cells while
committing six address columns.

Flock~\cite{flock} proves the compression relation with a batched R1CS over a Boolean witness. The witness is packed into $q_\pkd$, with $64$ bits per $\K$-element. $q_\pkd$ is the one aligned exception to the Jagged layout (\S\ref{sec:stacking}): it occupies the offset-zero prefix, so its ring-switch and fixed-slot claims keep their direct subcube formulas. The concrete ring switch of \S\ref{sec:ringswitch} maps its $\F_2$ claims to one $\E$-weighted claim on $q_\pkd$, which joins the batched opening (\S\ref{sec:e2e-unrolled}).

The eighteen $64$-bit words of $a,b,c,cv,m_0,m_1$ are contained in $q_\pkd$ and need no
separate value-column commitment. Each is a strided point evaluation of
$q_\pkd$ and is linked directly to its single-word memory interaction. The
R1CS matrices prove the compression relation for those supplied values.
\begin{itemize}
\item \textbf{Columns:} $\pc,\fp$; operands $o_{A,0},o_{A,1},o_{B,0},o_{B,1},o_{CV},o_C$ and metadata $m_0,m_1$; their six starting addresses; sixteen per-word memory counts; bytecode count $r_{\mathrm{bc}}$. The eighteen values live in $q_\pkd$.
\item \textbf{Constraints:} the six address relations $a_X=\fp o_X$. Flock proves $c=\mathrm{BLAKE3}(a,b;cv,m_0,m_1)$.
\item \textbf{Flushes:}
  \begin{itemize}
  \item \textbf{State:} \pull\ $\tup{\dsep{ST},\pc,\fp}$, \push\ $\tup{\dsep{ST},\gen\cdot\pc,\fp}$.
  \item \textbf{Bytecode:} \pull\ $\tup{\dsep{BC},\pc,r_{\mathrm{bc}},\opc{B3},o_{A,0},o_{A,1},o_{B,0},o_{B,1},o_{CV},o_C,m_0,m_1}$, \push\ the same at $\gen\cdot r_{\mathrm{bc}}$.
  \item \textbf{Memory:} sixteen reads: two at $a_{X,i},\gen a_{X,i}$ for each input chunk, and four at $\gen^j a_{CV}$ and $\gen^j a_C$ for $j\in\{0,1,2,3\}$, each paired with the count multiplied by $\gen$.
  \end{itemize}
\end{itemize}

\section{End-to-end protocol}\label{sec:e2e}

This section assembles the primitives into the executable protocol, beginning with the constants, encodings, and padding fixed in advance.

\subsection{Constants}\label{sec:e2e-const}

Fixed once, public, and shared by prover and verifier: the fields, the bus shape, and the tags naming interactions and instructions on it.

\paragraph{Fields, generator, and bus.}
\[
\K=\F_2[x]/(x^{64}+x^4+x^3+x+1),\qquad
\E=\K[y]/(y^3+y+1),\qquad
\gen=x,\qquad \operatorname{ord}(\gen)=2^{64}-1.
\]
Committed columns and machine words are $\K$-valued; challenges are $\E$-valued (\S\ref{sec:prelim}). Every bus tuple has width $m=12$ (the bytecode tuple); shorter tuples are zero-padded in their high coordinates (\S\ref{sec:m3}).

\paragraph{Domain separators.} Coordinate $0$ of every tuple is a domain separator, naming the interaction it belongs to (\S\ref{sec:m3}):
\[
  \dsep{ST}=1,\qquad \dsep{MEM}=\gen,\qquad \dsep{BC}=\gen^{2}.
\]

\paragraph{Opcodes.} Coordinate $3$ of every bytecode tuple is an opcode, naming the instruction stored at that program counter (\S\ref{sec:bytecode}):
\[
\begin{aligned}
  \opc{XOR}&=1,& \opc{MUL}&=\gen,& \opc{SET}&=\gen^{2},& \opc{DRF}&=\gen^{3},& \opc{JMP}&=\gen^{4},\\
  \opc{B3}&=\gen^{5},& \opc{XOR\_192}&=\gen^{6},& \opc{MUL\_192}&=\gen^{7},& \opc{DRF\_192}&=\gen^{8}.&&
\end{aligned}
\]

\subsection{Bytecode encoding}\label{sec:e2e-bc}

The program is public. Instruction $\pc$ fills coordinates $3$ through $11$ of its bytecode tuple $\tup{\dsep{BC},\pc,\cnt,\cdot}$. Coordinate $3$ holds the opcode; coordinates $4$ through $8$ hold five operand/flag slots; and coordinates $9$ through $11$ hold BLAKE3's output address and two metadata words (zero for every other instruction). Each address operand is a $\gen$-power and each immediate is a base-field word.
\begin{center}
\begin{tabular}{llllllllll}
\hline
Instruction & $3$ & $4$ & $5$ & $6$ & $7$ & $8$ & $9$ & $10$ & $11$\\
\hline
\texttt{XOR}           & $\opc{XOR}$ & $o_A$ & $o_B$ & $o_C$ & $0$ & $0$ & $0$ & $0$ & $0$\\
\texttt{MUL}           & $\opc{MUL}$ & $o_A$ & $o_B$ & $o_C$ & $0$ & $0$ & $0$ & $0$ & $0$\\
\texttt{XOR\_192}      & $\opc{XOR\_192}$ & $o_A$ & $o_B$ & $o_C$ & $0$ & $0$ & $0$ & $0$ & $0$\\
\texttt{MUL\_192}      & $\opc{MUL\_192}$ & $o_A$ & $o_B$ & $o_C$ & $0$ & $0$ & $0$ & $0$ & $0$\\
\texttt{SET\_CONSTANT} & $\opc{SET}$ & $o$   & $k$ & $0$ & $0$ & $0$ & $0$ & $0$ & $0$\\
\texttt{DEREF}         & $\opc{DRF}$ & $\alpha$ & $\beta$ & $\gamma$ & $f_{\pc}$ & $f_{\fp}$ & $0$ & $0$ & $0$\\
\texttt{DEREF\_192}    & $\opc{DRF\_192}$ & $\alpha$ & $\beta$ & $\gamma$ & $0$ & $0$ & $0$ & $0$ & $0$\\
\texttt{JUMP}          & $\opc{JMP}$ & $o_c$ & $o_d$ & $o_f$ & $0$ & $0$ & $0$ & $0$ & $0$\\
\texttt{BLAKE3}        & $\opc{B3}$  & $o_{A,0}$ & $o_{A,1}$ & $o_{B,0}$ & $o_{B,1}$ & $o_{CV}$ & $o_C$ & $m_0$ & $m_1$\\
\hline
\end{tabular}
\end{center}
So the program is nine public columns, whose MLEs the verifier forms itself. For \texttt{DEREF}, two operand slots are the store-mode flags; for \texttt{BLAKE3}, the slots hold six addresses and two metadata words. Writing $\mathit{instr}$ for the nine columns, the bytecode seed at $\pc$ pushes $\tup{\dsep{BC},\pc,1,\mathit{instr}}$ and an instruction table pulls it with its opcode hardcoded (\S\ref{sec:tables}); balance then forces each table's operands to match the program at its $\pc$.

The program is padded to a power of two (\S\ref{sec:idxcol}). A padding entry is never executed: no table pulls it, so its seed push and finalize pull (both at count $1$) cancel.

\subsection{Padding rows}\label{sec:e2e-pad}

Each table is padded to a power of two. Padding rows are a fixed default, not all zeros (a zero count would self-cancel a read, \S\ref{sec:memchan}): every column zero but the read counts $\gen^{0}=1$, satisfying every constraint. The prover announces each table's real row count and the verifier divides the default surplus out of the bus product (\S\ref{sec:gp}).

\subsection{Public input}\label{sec:e2e-pi}

The public input is the first four memory cells $\mem[0],\dots,\mem[3]\in\K$.
They span a two-variable subcube of the address cube. The verifier samples
$(r_0,r_1)\in\E^2$, interpolates the four public words on that subcube, and
adds the single claim
\[
  \widetilde M(r_0,r_1,0,\dots,0)=
  \operatorname{MLE}_{2}(\mem[0],\mem[1],\mem[2],\mem[3];r_0,r_1).
\]
The final stacked opening discharges this memory-column claim.

\subsection{The unrolled protocol}\label{sec:e2e-unrolled}

We lay out the whole protocol as one ordered sequence. Each reduction below adds evaluation claims on committed columns to a shared \emph{claim pool}; the final \textsc{Opening} phase proves all of them with a single PCS opening.

\paragraph{Setup and statement binding.} Fixed and shared in advance (\S\ref{sec:e2e-const} through \S\ref{sec:e2e-pi}): the fields and constants; the instance caps (\S\ref{sec:memchan}); the public program and its real instruction-prefix length; the initial and final state; and the four public words. Before any challenge, the transcript is seeded by the public input and an environment digest binding the bytecode, its real prefix, and the BLAKE3 circuit family. The prover then announces the used-memory prefix and the nine real opcode-row counts; the logical memory size is canonically $2^h$ for $h=\max(16,\lceil\log_2 h_{\rm mem}\rceil)$.

\paragraph{Commitment.}
\begin{enumerate}
\item \textbf{Prover.} Builds the $\K$-valued table columns, the single memory-word column $M$, the memory and bytecode finalize counts, and flock's packed witness $q_\pkd$ (\S\ref{sec:tab-blake3}). Extension operands contribute three value columns but only one committed starting-address column per operand; successor addresses are virtual. A program with no \texttt{BLAKE3} execution carries one padding instance so the proof shape remains uniform. Counts are committed $\gen$-powers (\S\ref{sec:memchan}); the public program and virtual increments $\gen\cdot\pc,\gen\cdot\cnt$ are not. The prover packs the real column prefixes into the Jagged witness $q$ and sends its dense commitment (\S\ref{sec:stacking}, \S\ref{sec:rs-dense}). The verifier checks the instance caps and derives both the Jagged boundaries and the power-of-two logical AIR and grand-product shapes (\S\ref{sec:leafstack}). Padding rows are fixed defaults (\S\ref{sec:e2e-pad}); the verifier restores their omitted MLE contributions during opening.
\end{enumerate}

\paragraph{Bus.}
\begin{enumerate}
\item \textbf{Verifier.} Sends the fingerprint challenge $\alpha$ and the multiset challenge $\gamma$.
\item \textbf{Prover.} Forms the push, pull, and \emph{count} leaf vectors (\S\ref{sec:leafstack}; the count leaves are the count columns themselves, \S\ref{sec:memchan}), stacked into blocks and padded with $1$s. Sends the count root $R_c$ and one bus root $R$, the push/pull product after the default surplus is removed (\S\ref{sec:gp}).
\item \textbf{Prover \& Verifier.} Run the one batched GKR (\S\ref{sec:gkr}) over all three trees. For the bus, the verifier multiplies $R$ by each side's public surplus and checks the side's product against it; one $R$ for both sides is the balance check (\S\ref{sec:gp}). For the counts it checks $R_c\neq0$ (no read self-cancels, \S\ref{sec:memchan}). The pass yields three leaf claims $\widetilde V^{t}_0(\zeta)$ at one shared $\zeta$.
\item \textbf{Prover \& Verifier.} Decompose each leaf claim (\S\ref{sec:leafstack}). The verifier forms the public part: the block selectors; the constant coordinates (separators, the per-table read opcodes, the boundary state $(\pc_0,\fp_0)$ and $(\pc_{\mathrm{final}},\fp_{\mathrm{final}})$); the index column (\S\ref{sec:idxcol}); and the public program columns of the bytecode seed and finalize, whose MLEs it forms directly. For the table-less blocks the prover supplies the committed-column evaluations (one $\E$-element each, \S\ref{sec:prelim}); a virtual coordinate such as $\gen\cdot\pc$ is read off its committed source, and these evaluations enter the claim pool. It sends nothing for the other blocks: what they owe the side is $\widetilde V_0(\zeta)$ less the part just formed, so the verifier forms that too, and the batch below settles it. The count channel decomposes likewise; none of its blocks is table-less, so its whole share falls to the batch.
\end{enumerate}

\paragraph{Local constraints.} Run once per table (\S\ref{sec:air}); the nine runs share no challenges.
\begin{enumerate}
\item \textbf{Verifier.} Sends the constraint-batch challenge $\eta$. Table $T_j$ takes a disjoint range of its powers, one per constraint; the bus forms absorbed above take three further powers, one per side and shared by every table. No point is drawn, the batch running at the bus point with $T_j$ tested at $\zeta_{<\tau_j}$.
\item \textbf{Prover \& Verifier.} Run the back-loaded batched zerocheck, an $n$-round sumcheck with $n=\max_j\tau_j$, a table of $2^{\tau}$ rows joining at round $n-\tau$. It starts from a target the verifier derives from the three leaf claims, and each round is four $\E$-elements, the round polynomial itself at four nodes. Its end needs every column of every table, \texttt{BLAKE3}'s eighteen value columns included, whose claims route to $q_{\mathrm{pkd}}$; each at that table's own reduced point $\rho_{<\tau_j}$, all of them prefixes of one $\rho$; these enter the pool, and they are also what the bus forms are read off.
\end{enumerate}

\paragraph{Public input.}
\begin{enumerate}
\item \textbf{Verifier.} Sends $(r_0,r_1)$, forms the two-variable interpolation of the four public words, and adds the single memory-column claim of \S\ref{sec:e2e-pi}.
\end{enumerate}

\paragraph{BLAKE3 validity.}
\begin{enumerate}
\item \textbf{Prover \& Verifier.} Reduce the executed BLAKE3 constraints to two evaluation claims on $q_\pkd$ (\S\ref{sec:tab-blake3}).
\item \textbf{Prover \& Verifier.} The ring switching of \S\ref{sec:ringswitch} reduces each bit-witness claim to a weighted-sum claim on the $q_\pkd$ region of the stack, with a transparent $\E$-valued weight the verifier can evaluate; both join the pool.
\end{enumerate}

\paragraph{Opening.}
\begin{enumerate}
\item \textbf{Prover \& Verifier.} Each pooled claim is one evaluation of a committed column at a point, its value already supplied.  The verifier removes any fixed public-padding contribution and, via the interval indicator of \S\ref{sec:stacking}, turns the real-prefix claim into a weighted sum $\sum_w W_j(w)\,\widehat q(w)=c_j$ over the dense Jagged witness.
\item \textbf{Verifier.} Sends the batching challenge $\lambda$.
\item \textbf{Prover \& Verifier.} Fold the $J$ claims into $W_\lambda=\sum_j\lambda^{j}W_j$ and $C_\lambda=\sum_j\lambda^{j}c_j$, and run the inner-product PCS opening on $\sum_w W_\lambda(w)\,\widehat q(w)=C_\lambda$ (\S\ref{sec:rs-dense}), the verifier evaluating $W_\lambda$ itself. This one opening discharges every pooled claim; there is no separate reduction sumcheck.
\item \textbf{Verifier.} Accepts iff the announced sizes passed the caps, the two roots agreed after the surplus division, the count root was nonzero, every sumcheck's rounds were consistent (the batch's against the target derived from the three leaf claims) and its final value matched the supplied evaluations, flock's reduction verified, the public-input claim held, and the PCS opening verified.
\end{enumerate}

\section{Recursive aggregation}\label{sec:recursion}

The implemented N-to-1 aggregation runs the verifier of \S\ref{sec:e2e-unrolled} \emph{inside} the VM: one compiled guest replays the transcripts, sumchecks, ring switches, and stacked Ligerito openings of $n_{\mathrm{rec}}\ge1$ ordinary VM proofs for a fixed inner program. The proofs may have different announced table and commitment sizes; the guest dispatches among the certified shapes without recompilation. Base-field bookkeeping uses ordinary machine words, extension verifier arithmetic uses three-word values and the explicit extension opcodes, and Fiat--Shamir hashing uses \texttt{BLAKE3}.

\subsection{Deferred evaluation claims}\label{sec:deferred-claims}

Twice during verification, a large \emph{fixed} multilinear must be evaluated at a random point:
\begin{itemize}
\item \textbf{the bytecode}: during leaf decomposition (\S\ref{sec:e2e-unrolled}, \emph{Bus}), the verifier evaluates the nine public program columns, padded to sixteen selector slots under four selector variables, in one pass over the whole program;
\item \textbf{the flock matrices}: the \texttt{BLAKE3} relation is proven by flock (\S\ref{sec:tab-blake3}), whose lincheck ends by evaluating $\widetilde A_0$ and $\widetilde B_0$, the per-block R1CS matrices: $28$-variable multilinears with $\approx 21\mathrm{M}$ nonzero Boolean entries between them.
\end{itemize}
Either evaluation is one pass over a fixed object. This takes milliseconds natively, but in-circuit it would dwarf the entire remaining verifier, in cycles and in bytecode size alike.

The aggregation guest therefore does not evaluate these fixed tables. It verifies every surrounding reduction and exports only the resulting points and claimed values.

There are three fixed polynomials: the stacked bytecode and the two matrices $A_0,B_0$. Verifying each ordinary sub-proof produces one claim on each, for $3n_{\mathrm{rec}}$ claims. A fresh Fiat--Shamir transcript absorbs all sub-statements and claim data, then batches them with two sumchecks: one for bytecode and one for the matrices, which share their $28$ variables. The guest reduces the batch to one bytecode evaluation and a common-point evaluation of each matrix. These three reduced claims, the inner proving-environment digest, and the sub-statements are hashed into the outer VM's four-word public input.

The public \texttt{RecursiveProof::verify} path first verifies that outer VM proof and then evaluates the three fixed polynomials natively at the reduced points. It rejects an empty batch and is the only complete acceptance path. Feeding a \texttt{RecursiveProof} into another aggregation layer (carrying its reduced claims forward rather than checking them natively) is not yet exposed by the current harness.

The matrix sumcheck runs over the $28$-variable polynomial
\[
  P(x) \;=\; A_0(x)\,W_A(x) \;+\; B_0(x)\,W_B(x),
  \qquad
  W_M(x) \;=\; \textstyle\sum_t \gamma_{M,t}\,\eq(r_t,x),
\]
and ends in the reduced claims $\widetilde A_0(r^\ast)$ and $\widetilde B_0(r^\ast)$, at a common point. The in-circuit cost is that of the sumcheck verifier: $28$ cheap rounds. The prover exploits sparsity to avoid the $2^{28}$-sized cube. During the first $14$ rounds (the row variables), each claim enters through the column contraction $m_t[i]=\sum_{j\in\mathrm{row}_i}\eq(r_t^{\mathrm{col}},j)$; for the last $14$ rounds, each matrix is contracted once along its rows, at the row point just fixed. Each contraction is a single addition-only pass over the nonzeros (the entries are Boolean), and all remaining work is dense of size $2^{14}$. Batching $t$ claims per matrix therefore costs $t+1$ passes over its nonzeros plus $O(t\cdot2^{14})$ multiplications, against the $O(2^{28})$ field operations and memory of a dense sumcheck. The bytecode sumcheck is the same, only smaller and with a dense polynomial.

\section{ISA programming}\label{sec:isa-programming}

How the zkDSL compiler (\texttt{crates/lean\_compiler}; surface language in \texttt{crates/lean\_compiler/zkDSL.md}) programs the nine-instruction ISA. Two themes run through every pattern: \emph{write-once memory is an assertion mechanism} (a second write of a cell is an equality check, an unwritten cell is prover-chosen), and \emph{indices live in the exponent} (a logical index $i$ is the element $\gen^{i}$, so incrementing, offsetting, and address formation are single field multiplications).

\subsection{Hints}\label{sec:prog-hints}

The prover's nondeterminism enters through write-once cells that are still unset when execution needs them. Some are explicit advice (a fresh allocation pointer $\gen^{\mathrm{base}}$, a witness stream, a computed-advice builtin); others are back-solved by an instruction whose remaining operands are known. In every case the resulting word is unconstrained except by the instructions and memory equalities that consume it, so a program must check all security-relevant advice.

\subsection{Division and inequality}\label{sec:prog-div-ne}

The zkDSL's single-slash division $q=a/b$ uses write-once back-solving: it emits \texttt{MUL} with $q$ as the one unset input and $a$ as the already-written output. Witness generation fills $q=a\,b^{-1}$, while the ordinary multiplication table proves $q b=a$. It costs one VM instruction. Division by zero is undefined. Double slash \texttt{//} remains compile-time integer floor division for sizes and indices; it is not a field operation.

The assertion \texttt{assert a != b} first computes $a+b$ with \texttt{XOR}. A conditional \texttt{JUMP} skips a poison path exactly when this value is nonzero. Equality falls through to a jump to $\gen^{-1}$, outside the committed bytecode cube, so the bytecode bus cannot balance any purported continuation. This proves inequality without an inverse hint.

\subsection{Functions}\label{sec:prog-functions}

A frame is fp-relative with layout $[\mathit{retpc},\mathit{retfp},\mathrm{args},\mathrm{rets},\mathrm{locals}]$. A call: the callee frame pointer is hinted into a fresh cell $\mathit{nfp}$; the arguments and the caller's $\fp$ are stored through $\mathit{nfp}$ with \texttt{DEREF}s (modes \texttt{cell} and \texttt{fp}); a \texttt{DEREF} in mode \texttt{pc} stores the return address $\gen^{2}\cdot\pc$ (the instruction after the transfer); one \texttt{JUMP} enters the callee. Return is a single \texttt{JUMP} to $\loc{1}$ with frame pointer $\loc{\gen}$. Cost: $n_{\mathrm{args}}+n_{\mathrm{rets}}+4$ instructions per call.

\subsection{Loops}\label{sec:prog-loops}

A counted loop keeps its counter in the exponent: it starts at $\gen^{\mathrm{lo}}$, advances by one multiplication $i\gets\gen\cdot i$, and exits on reaching $\gen^{\mathrm{hi}}$. The body is lowered as a tail-recursive helper function whose exit test \emph{is} the recursive call's \texttt{JUMP} condition (the \texttt{XOR} $i+\gen^{\mathrm{hi}}$ is nonzero exactly while the loop must continue), so an iteration costs one call plus two instructions, with no is-zero gadget. Loop state is threaded through captured arguments or a heap buffer.

\subsection{Conditionals}\label{sec:prog-conditionals}

\texttt{if}/\texttt{else} on a field equality is one \texttt{XOR} plus one conditional \texttt{JUMP}: the taken jump goes to whichever block the test should \emph{not} fall into ($=$ falls into \emph{then}, $\neq$ into \emph{else}), so no negation is ever computed. Intra-function jump targets are bytecode addresses $\gen^{\mathrm{entry}+i}$, backpatched at layout.

A taken \texttt{JUMP} reloads $\fp$ from a cell, and the ISA has no $\fp$-read; a branching function therefore materializes its own $\fp$ once: a \texttt{DEREF} in mode \texttt{fp} writes $\fp$ into a fresh one-cell heap buffer, a \texttt{DEREF} in mode \texttt{cell} copies it back into the frame (2 instructions, amortized; in \texttt{main}, $\fp=\gen^{0}=1$ is already a constant). Bindings made inside a branch are local to it; branches communicate through write-once cells: only one branch executes, so both may write the \emph{same} cell.

\subsection{Range checks}\label{sec:prog-range-checks}

The range check \emph{in the exponent} proves $x\in\{\gen^{0},\dots,\gen^{k-1}\}$, i.e.\ $\log_{\gen}x<k$, in 3 cycles (leanVM's \texttt{DEREF} trick \cite{leanvm-doc}, transported to $\gen$-powers):
\begin{enumerate}
\item \texttt{DEREF} through $x$: the dereferenced address $x\cdot\gen^{0}$ rides the memory interaction, which balances only for the $2^{h}$ addresses $\gen^{0},\dots,\gen^{2^{h}-1}$ (\S\ref{sec:memchan}): the bus itself proves $x=\gen^{e}$ with $e<2^{h}$;
\item \texttt{MUL} $x\cdot y$ into a write-once cell holding the constant $\gen^{k-1}$: the runner back-solves the complement $y=\gen^{k-1-e}$ (\S\ref{sec:prog-hints}), and the double write asserts $x\cdot y=\gen^{k-1}$;
\item \texttt{DEREF} through $y$: proves $y=\gen^{f}$, $f<2^{h}$.
\end{enumerate}
Then $e+f\equiv k-1\pmod{2^{64}-1}$ with $e,f<2^{h}\le2^{32}$, so $e+f<2^{33}<2^{64}-1$ and the congruence is an equality over the integers; a ``negative'' $k-1-e$ would wrap to $\approx2^{64}\gg2^{h}$, so $e\le k-1$, for every memory size the prover may announce, provided $k\le2^{16}$ (the minimum, \S\ref{sec:vm}). The two \texttt{DEREF} targets are deferred touches; the constant cell is one amortized \texttt{SET} per distinct bound per frame.

\subsection{Match statements}\label{sec:prog-match}

leanVM dispatches a switch with a single jump to $\pc=a+b\cdot x$: the case blocks, padded to a common length $b$, sit consecutively from base $a$, and its integer arithmetic forms the affine address in one instruction. That computation does not transport to this VM. Addresses here are $\gen$-powers: the additive offset $a+j$ is the cheap product $\gen^{a}\cdot x$, but the \emph{scaling} $b\cdot j$ is the power $x^{b}$ ($\log_2 b$ squarings per dispatch), with every block still padded to the longest. Instead, the aligned region collapses to fixed two-instruction \emph{trampoline slots} and the real blocks move out of line: matching on $x=\gen^{j}$, $j\in[0,n)$, the dispatch jumps to
\[
  d \;=\; \gen^{T}\cdot x^{2} \;=\; \gen^{T+2j},
\]
slot $j$ of a table at bytecode base $T$, where a slot is $\texttt{SET}\ c=\gen^{\mathrm{block}_j}$ then $\texttt{JUMP}\ c$. The case blocks sit anywhere, unaligned and of any length; the scaling cost is the single squaring $x^{2}$. Two jumps per dispatch, $\approx7$ cycles independent of $n$.

The slots are two instructions rather than one because a \texttt{JUMP} reads its destination from a \emph{cell}: a one-instruction slot would need its target cell pre-written, i.e.\ $n$ \texttt{SET}s executed before every dispatch. (Other layouts exist and trade exactly this, e.g.\ a memory-resident table of the $n$ block addresses, read with one \texttt{DEREF} and taken with a \emph{single} jump, paying the \texttt{SET}s as setup; worthwhile for many repeated small matches. The compiler currently emits only the trampoline.)

As in leanVM, the matched value must be known to lie in $[0,n)$: a hinted $x$ must be range-checked first (\S\ref{sec:prog-range-checks}); otherwise the dispatch lands at an attacker-chosen program counter.

\appendix

\section{Ring switching}\label{sec:ringswitch}

Flock's BLAKE3 argument produces claims about Boolean multilinears. The VM packs them into \(\K\)-words for commitment and opens the resulting claims over \(\E\). We use the following specialization of ring switching~\cite{cryptoeprint:2024/504}:
\[
  \F_2 \subset \K=\F_{2^{64}} \subset \E=\F_{2^{192}}.
\]

\paragraph{Packing.} Let \(L\) be the number of Boolean variables left after the six variables that select a bit within a \(64\)-bit word. Write the Boolean polynomial as
\[
  Q:\{0,\ldots,63\}\times\cube L\longrightarrow\F_2.
\]
Thus \(Q\) has \(L+6\) Boolean variables and \(64\cdot2^L\) values. In the polynomial basis \(1,x,\ldots,x^{63}\) of \(\K/\F_2\), pack each word as
\[
  Q_\pkd(y)=\sum_{i=0}^{63}Q(i,y)x^i\in\K,
  \qquad y\in\cube L.
\]
The PCS commits the \(2^L\) values of \(Q_\pkd\), stored in the \(q_\pkd\) region of the witness stack. This represents the Boolean witness with no embedding overhead.

\paragraph{Incoming claim.} Flock supplies \(a\in\E\), \(z\in\E^L\), and verifier-known weights \(p_i\in\E\) for \(0\le i<64\). A valid witness satisfies
\begin{equation}
  a=\sum_{i=0}^{63}p_i s_i,
  \qquad
  s_i=\sum_{y\in\cube L}\eq(z,y)Q(i,y).
  \label{eq:rs-incoming}
\end{equation}
The \(p_i\) are the Lagrange weights of Flock's univariate skip. The prover sends \(s_0,\ldots,s_{63}\in\E\), and the verifier checks the first equality. It remains to bind these values to the committed \(Q_\pkd\).

\paragraph{Coordinate transpose.} Let \(\vartheta\) be the generator of the cubic extension \(\E/\K\), and let \(e_{64j+i}=x^i\vartheta^j\), for \(0\le i<64\) and \(0\le j<3\), be the resulting basis of \(\E/\F_2\). Write
\[
  s_i=\sum_{w=0}^{191}a_{w,i}e_w,
  \qquad a_{w,i}\in\F_2,
  \qquad
  t_w=\sum_{i=0}^{63}a_{w,i}x^i\in\K.
\]
The map \((s_i)_{i<64}\mapsto(t_w)_{w<192}\) only transposes a \(192\times64\) bit matrix. Likewise, write
\[
  \eq(z,y)=\sum_{w=0}^{191}A_w(y)e_w,
  \qquad A_w(y)\in\F_2.
\]
Taking coordinates in the definition of \(s_i\) gives
\begin{equation}
  t_w
  =\sum_{i=0}^{63}x^i\sum_{y\in\cube L}A_w(y)Q(i,y)
  =\sum_{y\in\cube L}A_w(y)Q_\pkd(y),
  \qquad 0\le w<192.
  \label{eq:rs-coordinate-claim}
\end{equation}

\paragraph{Batching.} After absorbing all \(s_i\), the verifier samples independent \(f_0,\ldots,f_5\leftarrow\E\). Starting with \(v_0=v\), define
\[
  \begin{gathered}
    (d_0,\ldots,d_5)=(32,16,8,4,2,1),\qquad v_{j+1}=v_j+f_jv_j^{2^{d_j}},\\
    \Phi(v)=v_6,\qquad \lambda_w=\Phi(e_w).
  \end{gathered}
\]
Multiplying \eqref{eq:rs-coordinate-claim} by \(\lambda_w\) and summing gives
\begin{equation}
  T:=\sum_{w=0}^{191}\lambda_wt_w
  =\sum_{y\in\cube L}\underbrace{\Phi(\eq(z,y))}_{B_z(y)}Q_\pkd(y).
  \label{eq:rs-batched}
\end{equation}
This is directly a weighted opening of the committed \(Q_\pkd\): the target is \(T\in\E\), and \(B_z\) is a transparent \(\E\)-valued weight. No ring-switch sumcheck is needed. Multiple Flock claims share the same map; independent \(\E\)-scalars then fold their targets and weights into the global stacked opening.

\paragraph{Soundness.} Expanding the composition gives \(\Phi(v)=\sum_{k<64}C_k(f_0,\ldots,f_5)v^{2^k}\), with one distinct challenge monomial \(C_k\) for every \(k\). The generalized ring-switching analysis shows that any nonzero coordinate error makes at least one coefficient of \(\sum_k C_kW_k\) nonzero. Its total degree is at most \(2^{31}+2^{15}+2^7+2^3+2+1<2^{32}\), so sampling the challenges after the \(s_i\) bounds the batching error by less than \(2^{-160}\) before the PCS decoder's list-size factor. The full argument is in \href{https://github.com/leanEthereum/leanVM-b/blob/main/misc/ring-switching.pdf}{\texttt{ring-switching.pdf}}.

\paragraph{Succinct verification.} At a PCS query point \(r\in\E^L\), the verifier needs
\[
  \widehat B_z(r)=\sum_{y\in\cube L}\eq(r,y)\Phi(\eq(z,y))=\sum_{k<64}C_k\prod_{j=1}^{L}\left(1+r_j+z_j^{2^k}\right).
\]
The last equality follows by expanding \(\Phi\), factorizing the Boolean sum coordinate by coordinate, and using characteristic two. The six-stage recurrence generates the \(C_k\) once. The verifier then evaluates the expression in \(O(64L)\) field operations, without constructing the \(2^L\)-entry table. Applying \(\Phi\) to an ordinary field element is cheaper still: the composed form takes exactly \(63\) squarings and six multiplications.

\begin{thebibliography}{9}

\bibitem{m3}
Irreducible.
\newblock Multi-Multiset Matching (M3).
\newblock Binius documentation.
\newblock \url{https://www.binius.xyz/basics/binius-v0/m3/}

\bibitem{basefold}
Hadas Zeilberger, Binyi Chen, and Benjamin Fisch.
\newblock BaseFold: Efficient Field-Agnostic Polynomial Commitment Schemes from Foldable Codes.
\newblock Cryptology ePrint Archive, Paper 2023/1705, 2023.
\newblock \url{https://eprint.iacr.org/2023/1705}

\bibitem{cryptoeprint:2024/504}
Benjamin~E. Diamond and Jim Posen.
\newblock Polylogarithmic Proofs for Multilinears over Binary Towers.
\newblock Cryptology ePrint Archive, Paper 2024/504, 2024.
\newblock \url{https://doi.org/10.1007/978-3-032-25336-1_1}
\newblock \url{https://eprint.iacr.org/2024/504}

\bibitem{cryptoeprint:2024/108}
Angus Gruen.
\newblock Some Improvements for the PIOP for ZeroCheck.
\newblock Cryptology ePrint Archive, Paper 2024/108, 2024.
\newblock \url{https://eprint.iacr.org/2024/108}

\bibitem{cryptoeprint:2025/1187}
Andrija Novakovic and Guillermo Angeris.
\newblock Ligerito: A Small and Concretely Fast Polynomial Commitment Scheme.
\newblock Cryptology ePrint Archive, Paper 2025/1187, 2025.
\newblock \url{https://eprint.iacr.org/2025/1187}

\bibitem{whir}
Gal Arnon, Alessandro Chiesa, Giacomo Fenzi, and Eylon Yogev.
\newblock WHIR: Reed--Solomon Proximity Testing with Super-Fast Verification.
\newblock Cryptology ePrint Archive, Paper 2024/1586, 2024.
\newblock \url{https://eprint.iacr.org/2024/1586}
\bibitem{jagged-pcs}
Tamir Hemo, Kevin Jue, Eugene Rabinovich, Gyumin Roh, and Ron~D. Rothblum.
\newblock Jagged Polynomial Commitments (or: How to Stack Multilinears).
\newblock Manuscript, April 2026.

\bibitem{flock}
Succinct Labs.
\newblock flock: a batched-R1CS proving system for hash circuits over binary towers.
\newblock \url{https://github.com/succinctlabs/flock}

\bibitem{leanvm-doc}
leanEthereum.
\newblock leanVM: a minimal zkVM (\texttt{misc/minimal\_zkVM.tex}, \S Range checks, \S Switch statements).
\newblock \url{https://github.com/leanEthereum/leanVM}

\end{thebibliography}

\end{document}
